% Editorial bootstrap from the immutable source revision.
% This file is canonical and may be edited independently.

\part{Inverse Iterates and Their Taylor Obstructions}


% Source fragment: Inverse iterates; prefix ii.
\section{Statement of the problem and strengthened results}
\label{ii:sec:statement}

The bounded Fabius function is the unique smooth increasing bijection
$F\colon[0,1]\to[0,1]$ satisfying
\begin{equation}
 F(0)=0,\qquad F(1-x)=1-F(x),\qquad F'(x)=2F(2x)
 \quad(0\leq x\leq\half),
 \label{ii:eq:basic-F}
\end{equation}
with the usual constant continuation outside $[0,1]$ when the bounded totalization is
used.  Put
\begin{equation}
 I:=F^{-1},\qquad G_n:=F^{\comp n},\qquad K_n:=I^{\comp n},
 \qquad G_0=K_0=\id.
 \label{ii:eq:iterates}
\end{equation}
Associativity of composition gives the exact inverse identity
\begin{equation}
 K_n=G_n^{-1},\qquad G_n=K_n^{-1}.
 \label{ii:eq:iterate-inverse-identity}
\end{equation}
The inverse is $C^\infty$ on $(0,1)$, though it has singular endpoint slopes.

\begin{theorem}[Main theorem: inverse iterates are nowhere analytic]
\label{ii:thm:inverse-main}
For every integer $n\geq1$, the $n$-fold inverse iterate
\[
 K_n=I^{\comp n}
\]
is not real analytic at any point of $[0,1]$.
\end{theorem}

\begin{proof}
First let $y\in(0,1)$ and put $x=K_n(y)$.  The derivative formula
\eqref{ii:eq:inverse-iterate-derivative} gives $K_n'(y)>0$.  If $K_n$ were
analytic at $y$, the analytic inverse-function theorem would make its local
inverse $G_n$ analytic at $x$, contradicting \Cref{ii:thm:main}.  At $0$,
the quotient argument \eqref{ii:eq:Kn-no-derivative} shows that $K_n$ has no
finite right derivative, so it cannot be analytic even in the one-sided
endpoint sense.  Reflection, $K_n(1-y)=1-K_n(y)$, gives the same contradiction
at $1$.  Under the bounded totalization convention there is also a shorter
endpoint obstruction: $K_n$ is constant on the exterior side but strictly
increasing on the interior side, which is incompatible with an analytic germ.
\end{proof}

The proof yields a stronger statement about formal Taylor series.  For a smooth
function $h$ near $x$, write
\[
 R_h(x):=\rad\sum_{m\geq0}\frac{h^{(m)}(x)}{m!}z^m
\]
for the radius of convergence of its formal Taylor series; the series need not
represent $h$ even when this radius is positive.

\begin{theorem}[Co-countable zero-radius set for inverse iterates]
\label{ii:thm:inverse-strong}
For each fixed $n\geq1$ there exists a co-countable dense set
$Y_n\subset(0,1)$ such that
\begin{equation}
 R_{K_n}(y)=0,
 \qquad
 \limsup_{m\to\infty}
 \left(\frac{|K_n^{(m)}(y)|}{m!}\right)^{1/m}=+\infty
 \quad(y\in Y_n).
 \label{ii:eq:inverse-zero-radius}
\end{equation}
More precisely, one may take $Y_n=G_n(E_n)$, where the explicit co-countable dense set
$E_n$ is defined in \eqref{ii:eq:En} from forward orbit arithmetic.
\end{theorem}

\begin{proof}
Take the co-countable dense set $E_n$ of \eqref{ii:eq:En} and put
\begin{equation}
 Y_n:=G_n(E_n).
 \label{ii:eq:Yn}
\end{equation}
The homeomorphism $G_n:(0,1)\to(0,1)$ preserves density and
sends the countable complement of $E_n$ to the countable complement of
$Y_n$.  For $y=G_n(x)$ with $x\in E_n$, \Cref{ii:thm:strong} gives
$R_{G_n}(x)=0$, and the formal-reversion theorem
\Cref{ii:thm:formal-reversion} gives $R_{K_n}(y)=0$.  The displayed limsup is
then exactly the Cauchy--Hadamard formula.  The spine estimate establishing
the premise is proved in \Cref{ii:sec:proof-main}.
\end{proof}

The exceptional points are genuine and instructive.

\begin{theorem}[Affine but nonrepresenting center jet]
\label{ii:thm:center-jet}
For every $n\geq1$,
\begin{equation}
 K_n(\half)=\half,\qquad K_n'(\half)=2^{-n},\qquad
 K_n^{(m)}(\half)=0\quad(m\geq2).
 \label{ii:eq:center-jet}
\end{equation}
Thus the Taylor series of $K_n$ at $1/2$ is an affine polynomial of infinite radius,
but it does not agree with $K_n$ on any neighborhood.
\end{theorem}

\begin{proof}
The exact centre derivatives are
$F(\half)=\half$, $F'(\half)=2$, and $F^{(m)}(\half)=0$ for $m\ge2$.
Writing $\mathsf T_{\half}[G_n]$ for the formal Taylor series at the centre,
formal composition gives
$\mathsf T_{\half}[G_n](h)=\half+2^nh$.  Its unique formal compositional inverse
is $\half+2^{-n}h$, which proves all three derivative identities for $K_n$.
If that affine Taylor series represented $K_n$ on a neighbourhood, then
$K_n$ would be analytic there, contrary to \Cref{ii:thm:inverse-main}.
\end{proof}

The forward theorem needed for the proof is itself stronger than nowhere analyticity.

\begin{theorem}[Forward engine theorem]
\label{ii:thm:main}
For every integer $n\geq1$, $G_n=F^{\comp n}$ is not real analytic at any point of
$[0,1]$.
\end{theorem}

\begin{proof}
For the co-countable dense set $E_n$ in \eqref{ii:eq:En}, the finite spine
expansion \Cref{ii:thm:spine-expansion}, strict ratio unimodality
\Cref{ii:lem:prefix-unimodal}, the binary-transition lemma
\Cref{ii:lem:binary-transition}, and the defect bound
\Cref{ii:prop:defect-decay} give the dominant-spine estimate
\Cref{ii:prop:dominant-lower-bound}.  Hence \Cref{ii:thm:strong} gives zero
Taylor radius at every point of $E_n$.  Were $G_n$ analytic at any point of
$[0,1]$, its power series would represent it throughout a neighbourhood;
that neighbourhood meets the dense set $E_n$, a contradiction.  This also
covers the endpoints because every endpoint neighbourhood meets $(0,1)$.
\end{proof}

\begin{theorem}[Forward co-countable zero-radius set]
\label{ii:thm:strong}
Fix $n\geq1$.  There exists a co-countable dense set $E_n\subset(0,1)$ such that, for
every $x\in E_n$,
\begin{equation}
 \limsup_{m\to\infty}
 \left(\frac{|G_n^{(m)}(x)|}{m!}\right)^{1/m}=+\infty.
 \label{ii:eq:infinite-limsup}
\end{equation}
Consequently $R_{G_n}(x)=0$ for every $x\in E_n$.
\end{theorem}

\begin{proof}
The set $E_n$ in \eqref{ii:eq:En} is co-countable and dense because it removes
only one point, a finite tie set, and finitely many homeomorphic preimages of
the countable dyadic set.  At $x\in E_n$, the dominant-spine lower bound
\Cref{ii:prop:dominant-lower-bound} holds for infinitely many $m$:
\[
 |G_n^{(m)}(x)|\ge c(x,n)Q_mH_n(x)^m,
 \qquad Q_m=2^{m(m+1)/2}.
\]
Since $m!\le m^m$, taking $m$-th roots gives
\begin{equation}
 \left(\frac{|G_n^{(m)}(x)|}{m!}\right)^{1/m}
 \ge c(x,n)^{1/m}H_n(x)\frac{2^{(m+1)/2}}m\longrightarrow\infty
 \label{ii:eq:root-divergence}
\end{equation}
along that subsequence.  Cauchy--Hadamard now gives $R_{G_n}(x)=0$.
\end{proof}

The logical order of the report is deliberate.  We first establish the forward engine
by a derivative-spine argument.  We then transfer both nonanalyticity and zero Taylor
radius through local inversion.  The endpoint and center analyses are treated
separately because the inverse is not a local diffeomorphism at the endpoints and
because the center has a convergent but nonrepresenting Taylor series.

% Source fragment: Inverse iterates; prefix ii.
\section{Why a direct dyadic-jet induction fails}

It is useful to see the obstruction before introducing the cure.  For two smooth
functions $f$ and $h$,
\begin{equation}
 (f\comp h)^{(m)}(x)
 =\sum_{\pi\in\Pi_m}
 f^{(|\pi|)}(h(x))\prod_{B\in\pi}h^{(|B|)}(x),
 \label{ii:eq:faa-set-partitions}
\end{equation}
where $\Pi_m$ is the set of partitions of $\SetOf{1,\ldots,m}$.  Even if
$h^{(r)}(x)=0$ for all large $r$, partitions with many small blocks can survive.
Similarly, even if the high derivatives of $f$ vanish at $h(x)$, the term
$f'(h(x))h^{(m)}(x)$ survives.  There is no pointwise ``eventually zero jet''
induction for $F^{\comp n}$.

At the opposite extreme, a crude absolute-value bound on \eqref{ii:eq:faa-set-partitions}
introduces Bell numbers.  Since the number of partitions is superexponential on an
ordinary exponential scale, such a bound loses the delicate comparison with $Q_m$.
The correct normalization must exploit the quadratic exponent in $Q_m$, not merely
count partitions.  This leads to the defect below.

\section{The dyadic defect of a Fa\`a di Bruno partition}
\label{ii:sec:defect}

Let $\pi\in\Pi_m$ have $k=|\pi|$ blocks of sizes
$r_1,\ldots,r_k$, so $r_1+\cdots+r_k=m$.  Define
\begin{equation}
 D(\pi):=\binom m2-\binom k2-\sum_{i=1}^k\binom{r_i}{2}.
 \label{ii:eq:defect}
\end{equation}
An equivalent and more transparent formula is
\begin{equation}
 D(\pi)=\sum_{1\leq i<j\leq k}(r_ir_j-1).
 \label{ii:eq:defect-pairs}
\end{equation}
Every summand is nonnegative.  Therefore $D(\pi)=0$ only in the two extremal cases:
there is one block, or every block is a singleton.  These are exactly the two
unbranched chain-rule spines.

\paragraph{Lean crosswalk: the finite positive-list defect is complete.}
\begingroup\small\sloppy
The module \path{PartitionDefect.lean} has exactly three public definitions
and thirty-three public theorems; the following inventory is exhaustive.
Its list \(r=[r_1,\ldots,r_k]\) records positive block sizes, but it does not
construct a set partition of \(\{1,\ldots,m\}\).
\begin{itemize}[leftmargin=2em]
 \item The three definitions are \path{pairSum}, \path{blockPairDefect}, and
 \path{partitionDefect}.
 \item The unordered-pair and triangular-number engine consists of
 \path{pairSum_nil}, \path{pairSum_cons}, \path{pairSum_one},
 \path{pairSum_add}, \path{pairSum_congr}, \path{pairSum_map},
 \path{choose_add_two}, and \path{choose_list_sum_two}.
 \item The conventional defect formulas, excess decomposition, and fixed-block
 lower bound are
 \path{blockPairDefect_eq_mul_sub_one},
 \path{partitionDefect_eq_pairSum_mul_sub_one},
 \path{partitionDefect_nonneg},
 \path{choose_sum_two_eq_choose_length_add_sum_add_partitionDefect},
 \path{partitionDefect_eq_choose_sum_sub_choose_length_sub_sum},
 \path{length_le_sum_of_pos}, \path{sum_map_sub_one},
 \path{pairSum_add_eq}, \path{pairSum_map_add_eq},
 \path{partitionDefect_eq_linear_add_pairwise_excess},
 \path{partitionDefect_lower_bound}, and
 \path{partitionDefect_fixed_block_bound}.
 \item The zero and sharp-equality classifiers are
 \path{pairSum_eq_zero_iff_pairwise},
 \path{blockPairDefect_eq_zero_iff},
 \path{partitionDefect_eq_zero_iff},
 \path{sub_one_mul_sub_one_eq_zero_iff},
 \path{partitionDefect_eq_lower_bound_iff}, and
 \path{partitionDefect_fixed_block_eq_iff}.
 \item The first-positive-shell closure is
 \path{add_sub_one_le_mul_of_pos},
 \path{mul_eq_add_sub_one_iff},
 \path{firstShell_le_fixedBlockProduct},
 \path{fixedBlockProduct_eq_firstShell_iff},
 \path{firstShell_le_partitionDefect},
 \path{partitionDefect_eq_firstShell_iff}, and
 \path{partitionDefect_twoBlock_firstShell}.
\end{itemize}
For positive entries the module proves
\[
 D(r)=(k-1)\sum_i(r_i-1)
      +\sum_{i<j}(r_i-1)(r_j-1)
 \ge (k-1)(m-k).
\]
It proves \(D(r)=0\) exactly when \(k\le1\) or every entry is one; equality
in the fixed-\(m\), fixed-\(k\) bound exactly when every pair contains a
singleton (equivalently at most one entry exceeds one); and, under precisely
\(2\le k<m\), the first-shell bound \(D(r)\ge m-2\), with equality exactly at
\(k=2\) or \(k=m-1\) together with that same one-nonsingleton condition.  The
explicit profile \([m-1,1]\) attains the shell for \(m\ge3\).

This closure does \emph{not} formalize the set-partition family \(\Pi_m\),
the quadratic-scale factorization \eqref{ii:eq:Q-factorization}, the weighted
Bell sum \eqref{ii:eq:defect-decay}, its asymptotics, the two-spine reduction,
the finite spine expansion, or either the forward- or inverse-iterate theorem.
Those steps remain paper-only.
\endgroup

Let $u(\pi)$ be the number of singleton blocks and let
\begin{equation}
 s(\pi):=k-u(\pi)
 \label{ii:eq:s-nonsingleton}
\end{equation}
be the number of non-singleton blocks.

\begin{lemma}[Quadratic-scale factorization]
\label{ii:lem:Q-factorization}
For every partition $\pi\in\Pi_m$,
\begin{equation}
 \frac{Q_k\prod_{r_i\geq2}Q_{r_i}}{Q_m}
 =2^{s(\pi)-D(\pi)}.
 \label{ii:eq:Q-factorization}
\end{equation}
\end{lemma}

\begin{proof}
Using $\log_2Q_r=r(r+1)/2$ and omitting singleton blocks from the product,
\begin{align*}
 2\log_2\frac{Q_k\prod_{r_i\ge2}Q_{r_i}}{Q_m}
 &=k(k+1)+\sum_{r_i\ge2}r_i(r_i+1)-m(m+1)\\
 &=k^2+s(\pi)+\sum_{r_i\ge2}r_i^2-m^2.
\end{align*}
On the other hand, expanding \eqref{ii:eq:defect} gives
\[
 2D(\pi)=m^2-k^2+s(\pi)-\sum_{r_i\ge2}r_i^2.
\]
Subtracting proves \eqref{ii:eq:Q-factorization}.
\end{proof}

The next proposition is the combinatorial engine.  Its role is stronger than saying
that each nonextremal partition has positive defect: it shows that the sum over all
of them is still exponentially small.

\begin{proposition}[Weighted defect decay]
\label{ii:prop:defect-decay}
For each fixed $K>0$ there exist constants $C_K>0$ and $\rho_K\in(0,1)$ such that
\begin{equation}
 \sum_{\substack{\pi\in\Pi_m\\1<|\pi|<m}}
   (2K)^{s(\pi)}2^{-D(\pi)}
 \leq C_K\rho_K^m
 \qquad(m\geq2).
 \label{ii:eq:defect-decay}
\end{equation}
\end{proposition}

\begin{proof}
We give the full counting argument because this is the point at which a naive Bell
number estimate fails.

Write $k=|\pi|$, let
\begin{equation}
 e:=m-k
 \label{ii:eq:excess}
\end{equation}
be the total block excess over singletons, and let $r_{\max}=m-t$ be the largest block
size.  Convexity of $r\mapsto\binom r2$ shows that, for fixed $m$ and $k$, the defect
is minimized by one block of size $m-k+1$ and $k-1$ singleton blocks.  Hence
\begin{equation}
 D(\pi)\geq e(k-1)=e(m-e-1).
 \label{ii:eq:defect-e}
\end{equation}
The contribution between a largest block and all other blocks gives the second bound
\begin{equation}
 D(\pi)\geq (m-t)t-(k-1)\geq t(m-t-1).
 \label{ii:eq:defect-t}
\end{equation}

First consider partitions with $1\leq e\leq m/4$.  Their non-singleton blocks contain
at most $2e$ elements, because the number of such blocks is at most $e$.  Choosing
those elements and partitioning them gives at most
\begin{equation}
  (2e+1)m^{2e}(2e)^{2e}
 \label{ii:eq:count-near-singletons}
\end{equation}
possibilities.  Also $s(\pi)\leq e$, while \eqref{ii:eq:defect-e} gives
$D(\pi)\geq e(3m/4-1)$.  Uniformly for $1\leq e\leq m/4$, the base-two logarithm
of the counting and weight factor is at most
\[
 \log_2(2e+1)+4e\log_2m+e\log_2\max(1,2K).
\]
For fixed $K$ this is at most $em/4$ once $m$ is large enough, whereas the
negative defect exponent is at most $-e(3m/4-1)$.  Thus the $e$-th contribution
is at most $2^{-em/3}$ after enlarging the threshold once more, uniformly in
$e$.  The resulting geometric sum is $\BigO(2^{-m/3})$.

Next consider partitions with $1\leq t\leq m/4$.  Choose the $t$ elements outside a
largest block and partition them.  Up to harmless overcounting there are at most
$m^t t^t$ choices, and $s(\pi)\leq t+1$.  By \eqref{ii:eq:defect-t},
$D(\pi)\geq t(3m/4-1)$.  The same explicit estimate as above, now with $t$ in
place of $e$ (and one harmless extra factor $\max(1,2K)$), bounds the $t$-th
contribution by $2^{-tm/3}$ for all sufficiently large $m$, uniformly for
$1\leq t\leq m/4$.  Its sum is again exponentially small.

It remains to treat the middle region $e>m/4$ and $t>m/4$.  If $t\leq3m/4$, then
\eqref{ii:eq:defect-t} gives $D(\pi)\geq c_1m^2$ for all sufficiently large $m$ and an
absolute $c_1>0$.  Suppose instead that $t>3m/4$, so every block has size less than
$m/4$.  If $k>m/4$, then \eqref{ii:eq:defect-e} gives
$D(\pi)\geq c_2m^2$.  If $k\leq m/4$, then
\begin{align*}
 D(\pi)
 &=\sum_{i<j}r_ir_j-\binom k2\\
 &=\frac{m^2-\sum_i r_i^2}{2}-\binom k2
 >\frac{3m^2}{8}-\frac{m^2}{32}
 =\frac{11m^2}{32}.
\end{align*}
In the middle region, the total number of partitions is at most $m^m$ and the weight
$(2K)^{s(\pi)}$ is at most $\max(1,2K)^m$.  Both are overwhelmed by
$2^{-c m^2}$.  Combining the three regions proves \eqref{ii:eq:defect-decay}, after
increasing $C_K$ to cover the finitely many small values of $m$.
\end{proof}

\section{A two-spine composition lemma}
\label{ii:sec:two-spine}

The weighted defect estimate converts Fa\`a di Bruno's formula into an asymptotically
binary chain rule.

\begin{lemma}[Two-spine reduction]
\label{ii:lem:two-spine}
Let $h$ be smooth near $x$, let $f$ be smooth near $y=h(x)$, and put
$B=\AbsoluteValue{h'(x)}$.  Suppose
there are constants $H>0$ and $K\geq1$ such that
\begin{align}
 \AbsoluteValue{f^{(r)}(y)}&\leq Q_r &&(r\geq1),
 \label{ii:eq:f-envelope}\\
 \AbsoluteValue{h^{(r)}(x)}&\leq KQ_rH^r &&(r\geq2).
 \label{ii:eq:h-envelope}
\end{align}
Then, with $L=\max(B,H)$,
\begin{equation}
 (f\comp h)^{(m)}(x)
 =f'(y)h^{(m)}(x)+f^{(m)}(y)h'(x)^m
  +\LittleO(Q_mL^m)
 \quad(m\to\infty).
 \label{ii:eq:two-spine}
\end{equation}
\end{lemma}

\begin{proof}
In \eqref{ii:eq:faa-set-partitions}, remove the one-block partition and the all-singleton
partition.  They give exactly the two displayed terms.  For any remaining partition
$\pi$ with $k$ blocks, $u$ singleton blocks, and $s=k-u$ non-singleton blocks, the
corresponding term has absolute value at most
\[
 Q_k B^u\prod_{r_i\ge2}(KQ_{r_i}H^{r_i})
 \leq Q_mL^m(2K)^s2^{-D(\pi)},
\]
where \cref{ii:lem:Q-factorization} was used in the last step.  Summing and applying
\cref{ii:prop:defect-decay} bounds the entire remainder by
$C_KQ_mL^m\rho_K^m=\LittleO(Q_mL^m)$.
\end{proof}

\begin{remark}[Why this lemma is composition-stable]
The estimate does not require analyticity, Gevrey bounds of factorial type, or a
closed formula for every Bell polynomial.  It uses only the exact quadratic dyadic
scale $Q_r$ and an ordinary exponential envelope $H^r$ for the inner derivatives.
The conclusion has the same form, so it can be iterated a fixed number of times.
\end{remark}

\section[The orbit-spine asymptotic for compositional iterates]{The orbit-spine asymptotic for $F^{\comp n}$}
\label{ii:sec:spines}

Fix $x\in(0,1)$ and $n\geq1$.  Define the finite orbit
\begin{equation}
  x_j:=G_j(x)=F^{\comp j}(x),\qquad 0\leq j\leq n-1,
  \label{ii:eq:orbit}
\end{equation}
and the local slopes
\begin{equation}
  a_j:=F'(x_j)>0.
  \label{ii:eq:slopes}
\end{equation}
The prefix and suffix weights are
\begin{align}
 B_0&:=1,
 &B_j&:=\prod_{r=0}^{j-1}a_r=G_j'(x)\quad(j\geq1),
 \label{ii:eq:prefix}\\
 A_{j,n}&:=\prod_{r=j+1}^{n-1}a_r,
 &&0\leq j<n,
 \label{ii:eq:suffix}
\end{align}
where empty products equal $1$.  Finally set
\begin{equation}
 H_n(x):=\max_{0\leq j<n}B_j.
 \label{ii:eq:Hn}
\end{equation}

\begin{theorem}[Finite spine expansion]
\label{ii:thm:spine-expansion}
For each fixed $n\geq1$ and $x\in(0,1)$,
\begin{equation}
 G_n^{(m)}(x)
 =Q_m\left[
   \sum_{j=0}^{n-1}A_{j,n}B_j^m\ThetaF(2^m x_j)
   +\LittleO\bigl(H_n(x)^m\bigr)
 \right]
 \qquad(m\to\infty).
 \label{ii:eq:spine-expansion}
\end{equation}
\end{theorem}

\begin{proof}
We induct on $n$.  For $n=1$, $x_0=x$, $A_{0,1}=B_0=H_1=1$, and
\eqref{ii:eq:spine-expansion} is exactly \eqref{ii:eq:derivative-atlas}.

Assume the formula for $G_j$.  It immediately supplies a global-in-order envelope at
the fixed point $x$: there is $K\geq1$ such that
\begin{equation}
 \AbsoluteValue{G_j^{(r)}(x)}\leq KQ_rH_j(x)^r\qquad(r\geq2).
 \label{ii:eq:Gj-envelope}
\end{equation}
Indeed the asymptotic gives the estimate for all sufficiently large $r$, and $K$ can
be enlarged to absorb the finitely many remaining orders.

Apply \cref{ii:lem:two-spine} with $f=F$ and $h=G_j$.  Since
$h(x)=x_j$, $h'(x)=B_j$, $f'(x_j)=a_j$, and
$f^{(m)}(x_j)=Q_m\ThetaF(2^m x_j)$, we obtain
\begin{equation}
 G_{j+1}^{(m)}(x)
 =a_jG_j^{(m)}(x)
  +Q_mB_j^m\ThetaF(2^m x_j)
  +\LittleO\bigl(Q_mH_{j+1}(x)^m\bigr),
 \label{ii:eq:spine-induction-step}
\end{equation}
where $H_{j+1}=\max(H_j,B_j)$.  Multiplication by $a_j$ changes every old suffix
factor $A_{i,j}$ into $A_{i,j+1}$, while the new term has suffix factor
$A_{j,j+1}=1$.  This is precisely \eqref{ii:eq:spine-expansion} for $j+1$.
\end{proof}

\begin{remark}[Combinatorial interpretation]
Each summand in \eqref{ii:eq:spine-expansion} corresponds to a chain-rule tree with one
high-valence vertex, located at orbit level $j$, and unary vertices everywhere else.
All trees with two or more genuinely branching vertices are swallowed by the
$\LittleO(H_n^m)$ remainder.  This is why ``spine'' is an apt name.
\end{remark}

\section{Strict unimodality of the spine weights}
\label{ii:sec:unimodality}

The finite sum in \eqref{ii:eq:spine-expansion} could still cancel if several $B_j$ have
the same maximal exponential size.  The orbit dynamics show that this is exceptional.

\begin{lemma}[Decreasing slope sequence]
\label{ii:lem:decreasing-slopes}
Let $x\in(0,1)\setminus\SetOf{\half}$.  Then the sequence
\begin{equation}
 a_0,a_1,a_2,\ldots
 \label{ii:eq:a-decreasing}
\end{equation}
is strictly decreasing.
\end{lemma}

\begin{proof}
If $x<\half$, then \eqref{ii:eq:left-orbit} says $x_{j+1}<x_j<\half$.  Since $F'$ is
strictly increasing on the left half, $a_{j+1}=F'(x_{j+1})<F'(x_j)=a_j$.
If $x>\half$, then $x_{j+1}>x_j>\half$ and $F'$ is strictly decreasing on the right
half, giving the same conclusion.
\end{proof}

\begin{lemma}[Unique or adjacent double maximum]
\label{ii:lem:prefix-unimodal}
For $x\in(0,1)\setminus\SetOf{\half}$, the finite positive sequence
$B_0,\ldots,B_{n-1}$ is strictly log-concave in the ratio sense:
\begin{equation}
 \frac{B_{j+1}}{B_j}=a_j
 \quad\text{is strictly decreasing in }j.
 \label{ii:eq:ratio-decreasing}
\end{equation}
Hence its maximum is unique unless $a_k=1$ for some $0\leq k<n-1$; in that case the
only maximizers are the adjacent pair $B_k=B_{k+1}$.
\end{lemma}

\begin{proof}
Equation \eqref{ii:eq:ratio-decreasing} follows from \eqref{ii:eq:prefix} and
\cref{ii:lem:decreasing-slopes}.  A positive finite sequence whose successive ratios
strictly decrease first rises while the ratio exceeds $1$, possibly has one flat step
when the ratio equals $1$, and then falls.  The conclusion follows.
\end{proof}

Because $F'$ rises strictly from $0$ to $2$ on $[0,\half]$, there is a unique
$u\in(0,\half)$ with $F'(u)=1$.  Symmetry gives the only other solution $1-u$.
For a fixed $n$, the tie set is therefore
\begin{equation}
 T_n:=\bigcup_{\substack{k\geq0\\k+1<n}}
 G_k^{-1}\!\left(\SetOf{u,1-u}\right).
 \label{ii:eq:tie-set}
\end{equation}
Thus $T_1=\varnothing$.  Each $G_k$ is a homeomorphism of $[0,1]$, so in
general $T_n$ is finite, with at most $2(n-1)$ points.

\section{Binary transitions force recurrent large amplitudes}
\label{ii:sec:binary}

Let $\D_2$ denote the dyadic rationals in $(0,1)$.

\begin{lemma}[Binary-transition lemma]
\label{ii:lem:binary-transition}
If $y\in(0,1)\setminus\D_2$, then there are infinitely many $m\geq1$ such that
\begin{equation}
 \AbsoluteValue{\ThetaF(2^m y)}\geq\half.
 \label{ii:eq:theta-recurrence}
\end{equation}
\end{lemma}

\begin{proof}
Choose the nonterminating binary expansion
\[
 y=0.b_1b_2b_3\cdots\quad(b_j\in\SetOf{0,1}).
\]
Since $y$ is not dyadic, its digits are not eventually constant: eventual zeros give
a terminating expansion, while eventual ones give the alternative expansion of a
dyadic rational.  Hence there are infinitely many transitions
$b_m\neq b_{m+1}$.

Modulo $2$, the number $2^m y$ is
\[
 b_m+0.b_{m+1}b_{m+2}\cdots.
\]
At a transition this lies in $[\half,\tfrac32]$.  By the block formula
\eqref{ii:eq:theta-block}, the centered argument of $\RvachevUp$ then lies in
$[-\half,\half]$.  Since $\RvachevUp$ is even and decreasing on $[0,1]$,
\eqref{ii:eq:up-half} gives
\[
 \AbsoluteValue{\ThetaF(2^m y)}
 =\RvachevUp(\text{centered argument})\geq\RvachevUp(\half)=\half.
\]
There are infinitely many such $m$.
\end{proof}

\begin{remark}
The lemma uses only adjacent binary digit changes, not equidistribution or normality.
It therefore applies to every non-dyadic point, including highly non-generic automatic,
Sturmian, or sparse-digit numbers.
\end{remark}

\section{Proof of the zero-radius theorem}
\label{ii:sec:proof-main}

For fixed $n$, define
\begin{equation}
 E_n:=(0,1)\setminus\left(
 \SetOf{\half}\cup T_n\cup
 \bigcup_{j=0}^{n-1}G_j^{-1}(\D_2)
 \right).
 \label{ii:eq:En}
\end{equation}
Every $G_j$ is a homeomorphism, so each preimage $G_j^{-1}(\D_2)$ is countable.
The tie set $T_n$ is finite.  Thus $E_n$ is co-countable and dense in $(0,1)$.

\begin{proposition}[Dominant-spine lower bound]
\label{ii:prop:dominant-lower-bound}
Let $x\in E_n$, and let $j_*$ be the unique index for which
$B_{j_*}=H_n(x)$.  There are infinitely many $m$ such that
\begin{equation}
 \AbsoluteValue{G_n^{(m)}(x)}
 \geq c(x,n)Q_mH_n(x)^m
 \label{ii:eq:dominant-lower-bound}
\end{equation}
with a constant $c(x,n)>0$ independent of $m$.
\end{proposition}

\begin{proof}
Because $x\in E_n$, the orbit point $x_{j_*}$ is not dyadic.  By
\cref{ii:lem:binary-transition}, there is an infinite set $M\subset\NaturalNumbers$ such that
\begin{equation}
 \AbsoluteValue{\ThetaF(2^m x_{j_*})}\geq\half
 \qquad(m\in M).
 \label{ii:eq:dominant-theta}
\end{equation}
For every $j\neq j_*$, put $r_j=B_j/H_n(x)<1$.  Divide
\eqref{ii:eq:spine-expansion} by $Q_mH_n(x)^m$.  Along $M$, the dominant summand has
absolute value at least $A_{j_*,n}/2$, while the sum of all other spines is bounded by
\[
 \sum_{j\neq j_*}A_{j,n}r_j^m\longrightarrow0.
\]
The normalized remainder also tends to zero.  Therefore, for all sufficiently large
$m\in M$, the normalized absolute value is at least $A_{j_*,n}/4$.  We may take
$c(x,n)=A_{j_*,n}/4>0$.
\end{proof}

This proposition is the last engine needed by the unique proofs of
\cref{ii:thm:strong,ii:thm:main} given with their statements.  Namely, taking
$m$th roots yields \eqref{ii:eq:root-divergence}; Cauchy--Hadamard gives zero
Taylor radius on the dense set $E_n$, and an analytic neighbourhood cannot
meet such a point.

\section{The inverse-iterate theorem}
\label{ii:sec:inverse-transfer}

\subsection{Interior diffeomorphism structure}

Let $I=F^{-1}$.  Since $F'(x)>0$ for $0<x<1$, every $G_n$ is a strictly increasing
$C^\infty$ diffeomorphism of $(0,1)$ onto itself, and its inverse is $K_n$.  For
$y\in(0,1)$ and $x=K_n(y)$,
\begin{equation}
 K_n'(y)=\frac{1}{G_n'(x)}
 =\prod_{j=0}^{n-1}\frac{1}{F'(F^{\comp j}(x))}>0.
 \label{ii:eq:inverse-iterate-derivative}
\end{equation}
Equivalently, writing $y_j=I^{\comp j}(y)$,
\begin{equation}
 K_n'(y)=\prod_{j=0}^{n-1} I'(y_j).
 \label{ii:eq:inverse-product-slope}
\end{equation}

The nonzero derivative in \eqref{ii:eq:inverse-iterate-derivative} is the
interior mechanism used in the unique proof of \cref{ii:thm:inverse-main}:
an analytic germ for $K_n$ would have an analytic local inverse, contrary to
\cref{ii:thm:main}.  That argument proves local nonrepresentability, but it
does not say whether the formal Taylor series of $K_n$ converges to some
unrelated analytic germ.  Formal reversion settles that stronger question.

\subsection{Formal-radius invariance under local inversion}

\begin{theorem}[Positive convergence radius is invariant under formal reversion]
\label{ii:thm:formal-reversion}
Let $g$ and $k$ be $C^\infty$ local inverses near $x_0$ and $y_0=g(x_0)$, with
$g'(x_0)\ne0$.  Then
\begin{equation}
 R_g(x_0)>0\quad\Longleftrightarrow\quad R_k(y_0)>0.
 \label{ii:eq:radius-invariance}
\end{equation}
Consequently $R_g(x_0)=0$ if and only if $R_k(y_0)=0$.
\end{theorem}

\begin{proof}
Translate the centers and write the formal Taylor series as
\[
 \mathsf T_g(z)=a_1z+a_2z^2+\cdots,
 \qquad
 \mathsf T_k(w)=b_1w+b_2w^2+\cdots,
 \qquad a_1b_1=1.
\]
Differentiating the smooth identities $g\comp k=\id$ and $k\comp g=\id$ to every
finite order gives the formal composition identities
\begin{equation}
 \mathsf T_g\comp\mathsf T_k=w,
 \qquad
 \mathsf T_k\comp\mathsf T_g=z.
 \label{ii:eq:formal-inverse-identities}
\end{equation}
Thus the two formal Taylor series are the unique formal compositional inverses
of one another.

Assume $R_k(y_0)>0$.  Then $\mathsf T_k$ defines an analytic germ with nonzero linear
coefficient.  By the analytic inverse-function theorem it has a convergent analytic
inverse germ.  The Taylor series of that inverse is the unique formal inverse of
$\mathsf T_k$, hence equals $\mathsf T_g$.  Therefore $R_g(x_0)>0$.  Interchanging $g$
and $k$ proves the converse.
\end{proof}

\begin{corollary}[Transport of the zero-versus-positive Taylor-radius locus]
\label{ii:cor:radius-transport}
For every $n\geq1$ and $x\in(0,1)$,
\begin{equation}
 R_{K_n}(G_n(x))=0\quad\Longleftrightarrow\quad R_{G_n}(x)=0.
 \label{ii:eq:radius-transport}
\end{equation}
The same equivalence holds with ``$=0$'' replaced by ``$>0$''.
\end{corollary}

\begin{proof}
At $x$ and $G_n(x)$, the smooth germs of $G_n$ and $K_n$ are mutually
inverse and $G_n'(x)>0$.  Applying \cref{ii:thm:formal-reversion} gives the
equivalence for positive radius.  A power-series radius is nonnegative, so
negating that equivalence gives the displayed zero-radius equivalence.
\end{proof}

Applied to the set $Y_n$ of \eqref{ii:eq:Yn}, the radius-transport corollary
is precisely the formal-series step in the unique proof of
\cref{ii:thm:inverse-strong}.  That theorem also gives an independent route
to interior nowhere analyticity: an analytic point has
an open analytic neighborhood, but every nonempty open interval meets the dense set
$Y_n$ on which the Taylor radius is zero.

\subsection{Endpoint obstruction and iterated inverse asymptotics}
\label{ii:sec:endpoint-inverse}

There is an elementary endpoint obstruction even before using the sharp asymptotic.
Since $F'(0)=0$,
\begin{equation}
 \frac{F(x)}{x}\longrightarrow0\qquad(x\downarrow0).
 \label{ii:eq:F-over-x-zero}
\end{equation}
Writing $y=F(x)$ gives
\begin{equation}
 \frac{I(y)}{y}=\frac{x}{F(x)}\longrightarrow+\infty.
 \label{ii:eq:I-over-y-infinite}
\end{equation}
Moreover $F(x)<x$ on $(0,1/2)$, so $I(y)>y$ there and therefore
$K_n(y)\geq I(y)$.  Hence
\begin{equation}
 \frac{K_n(y)}{y}\longrightarrow+\infty.
 \label{ii:eq:Kn-no-derivative}
\end{equation}
Thus $K_n$ has no finite right derivative at $0$ and cannot be analytic there.
The symmetry $I(1-y)=1-I(y)$, inherited by every iterate, gives the same conclusion at
$1$.  This is the endpoint ingredient invoked in the unique proof of
\cref{ii:thm:inverse-main}; it works under either the one-sided or totalized
endpoint convention.

The repository's formally backed endpoint theorem yields substantially more.  Put
$L=\log2$.  Its quadratic logarithmic law is
\begin{equation}
 -\log F(x)\sim\frac{(\log(1/x))^2}{2L}
 \qquad(x\downarrow0).
 \label{ii:eq:quadratic-endpoint-law}
\end{equation}

\begin{theorem}[Iterated inverse endpoint scale]
\label{ii:thm:iterated-endpoint}
For each fixed $n\geq1$, as $y\downarrow0$,
\begin{equation}
 -\log K_n(y)
 \sim c_n\bigl(\log(1/y)\bigr)^{2^{-n}},
 \qquad
 c_n=(2\log2)^{1-2^{-n}}.
 \label{ii:eq:iterated-inverse-endpoint}
\end{equation}
Consequently, for every $\alpha>0$,
\begin{equation}
 \frac{K_n(y)}{y^\alpha}\longrightarrow+\infty.
 \label{ii:eq:not-holder}
\end{equation}
Thus $K_n$ is not H\"older continuous of any positive order at either endpoint.
\end{theorem}

\begin{proof}
For $y\downarrow0$, let $T_0=\log(1/y)$ and recursively put
$T_j=\log(1/K_j(y))$.  If $u=I(v)$, then $v=F(u)$, so
\eqref{ii:eq:quadratic-endpoint-law} gives
\begin{equation}
 \log(1/u)\sim\sqrt{2L\log(1/v)}.
 \label{ii:eq:one-inverse-log}
\end{equation}
Hence $T_{j+1}\sim\sqrt{2LT_j}$.  Induction gives
\[
 T_n\sim c_n T_0^{2^{-n}},
 \qquad c_{n+1}=\sqrt{2Lc_n},\quad c_0=1.
\]
Solving the constant recursion yields
$c_n=(2L)^{1/2+1/4+\cdots+1/2^n}=(2L)^{1-2^{-n}}$, proving
\eqref{ii:eq:iterated-inverse-endpoint}.  Since $T_n=o(T_0)$,
\[
 \log\frac{K_n(y)}{y^\alpha}=\alpha T_0-T_n\longrightarrow+\infty,
\]
which proves \eqref{ii:eq:not-holder}.  Reflection gives the endpoint $1$.
\end{proof}

The sharper repository expansion inserts a lower-branch Lambert--$W$ phase and a tiny
period-one ripple.  Iterating that full expansion should produce a hierarchy of nested
Lambert phases; the coarse law above is the exact leading scale and is all that is
needed for the regularity theorem.

\subsection{The fixed center: a convergent Taylor series that fails to represent}

The proof of \Cref{ii:thm:center-jet} was placed with its statement.  Its
conclusion is worth emphasizing: the entire nonlinear remainder is flat at
the centre but nonzero on every punctured neighbourhood.

\subsection{No local elementary representation}

The repository's non-elementarity development proves that every function in its
formal elementary real-variable class has a dense open analytic locus.

\begin{corollary}[Local non-elementarity]
\label{ii:cor:non-elementary}
Fix $n\geq1$.  No elementary function in the repository's class can agree with
$K_n=I^{\comp n}$ on a subset of $(0,1)$ having nonempty interior.  The same statement
holds for $G_n=F^{\comp n}$.
\end{corollary}

\begin{proof}
Agreement on an open interval with a function having a dense analytic locus would make
the corresponding iterate analytic at some point of that interval, contrary to
\cref{ii:thm:inverse-main,ii:thm:main}.
\end{proof}

\subsection{What formal reversion does and does not preserve}

The forward lower bound on $E_n$ has the explicit superfactorial scale
\begin{equation}
 |G_n^{(m)}(x)|\geq cQ_mH_n(x)^m
 \quad\text{for infinitely many }m,
 \qquad Q_m=2^{m(m+1)/2}.
 \label{ii:eq:forward-growth-restated}
\end{equation}
Formal reversion proves zero radius for the inverse Taylor series, but by itself it does
not identify a comparable pointwise leading term for $K_n^{(m)}$.  The inverse Bell
polynomials contain additional cancellation.  A direct inverse-spine asymptotic is
therefore a genuine next problem, not an automatic corollary.  The numerical section
suggests that the inverse Taylor roots grow even faster at the chosen diagnostic point,
but no asymptotic claim is based on those finite-order data.

% Source fragment: Inverse iterates; prefix ii.
\section{Conjectures and research directions}
\label{ii:sec:future}

The main theorem is complete, but it leaves a structured exceptional countable set on
which the formal Taylor series may converge without representing the function.  The
formal-radius transport theorem shows that the inverse classification is exactly the
homeomorphic image of the forward positive-radius locus.  Of the five submitted
conjecture environments, only the direct inverse-spine statement remains a manuscript
conjecture.  Three are retained as explicit warning records: one has a false inverse
clause, one is already proved from exact quarter-point facts, and one duplicates the
corrected forward report's sole live conjecture.  The former nested-Lambert statement
is demoted below to a precise research problem because its proposed coefficient scale,
phase recursion, and uniform remainder had not been specified.  Neither the surviving
conjecture nor that problem has an exact proved Lean counterpart.

\begin{warning}[Former Conjecture 14.1 has a false inverse clause]
\label{ii:warn:trichotomy}
The submitted statement first proposed that the positive-radius locus of $G_n$ is
precisely its eventually-zero-jet locus, then claimed that
\cref{ii:thm:formal-reversion} gives the same eventual-zero classification for $K_n$.
The second clause is false: formal reversion preserves convergence, not polynomiality.
At $x=1/4$ the exact forward Taylor germ is
\[
 \mathsf T_{1/4}[F](z)=\frac5{72}+z+4z^2,
\]
because $F'(1/4)=1$, $F''(1/4)=8$, and all higher derivatives vanish.  At
$y=F(1/4)=5/72$, the inverse Taylor germ therefore solves $w=z+4z^2$ and is
\[
 z=\frac{\sqrt{1+16w}-1}{8},
\]
an infinite convergent series.  Thus $K_1=I$ has positive Taylor radius at $5/72$
without an eventually-zero jet.  The forward-only assertion remains a plausible open
classification problem, especially for $n\geq2$, but no inverse eventual-zero
equivalence is asserted here.
\end{warning}

The center is exceptional in both directions because its forward and inverse formal
series are affine.  At noncentral dyadic points for $n=1$, however, finite nonlinear
forward Taylor polynomials revert to infinite convergent inverse series, exactly as the
quarter-point example shows.

\begin{warning}[Former Conjecture 14.2 is already discharged]
\label{ii:warn:ties}
The submitted tie-cancellation statement is not a live conjecture.  If
$x\in T_n$, the tied adjacent weights satisfy $a_k=1$.  Strict monotonicity of $F'$,
together with $F'(1/4)=F'(3/4)=1$, forces
$x_k\in\{1/4,3/4\}$.  Hence
$x_{k+1}=F(x_k)\in\{5/72,67/72\}$.  The first orbit point is dyadic, so
$\ThetaF(2^m x_k)=0$ for every sufficiently large $m$ by
\eqref{ii:eq:dyadic-jet}; the second is non-dyadic, so
\cref{ii:lem:binary-transition} supplies infinitely many $m$ with
$\AbsoluteValue{\ThetaF(2^m x_{k+1})}\geq1/2$.  Since the fixed suffix coefficient
$A_{k+1,n}$ is positive, along that subsequence the surviving tied spine term has
absolute value at least $(A_{k+1,n}/2)H_n(x)^m$.  Thus persistent cancellation is
already excluded by exact quarter-point facts and the binary-transition lemma; no
additional Thue--Morse correlation conjecture remains.
\end{warning}

\begin{conjecture}[Direct inverse-spine asymptotic]
\label{ii:conj:inverse-spine}
For $y=G_n(x)$ away from an explicit countable exceptional set, the derivatives of
$K_n$ admit an asymptotic expansion on the same universal quadratic scale $Q_m$,
\[
 K_n^{(m)}(y)=Q_m\left[
   \sum_{j=0}^{n-1}\widetilde A_{j,n}(x)
   \widetilde B_j(x)^m\Theta(2^m x_j)
   +o(\widetilde H_n(x)^m)
 \right],
\]
with computable inverse prefix weights $\widetilde B_j$ obtained from the forward
slopes and with a unique dominant term off finite tie surfaces.
\end{conjecture}

The implicit identity $G_n\comp K_n=\id$ and the same partition defect provide a
plausible route, but one must control the recursive inverse Bell polynomials without
assuming the desired envelope.  A proof would upgrade abstract zero-radius transport
to explicit derivative lower bounds for the inverse iterates.  The finite-order replay
does not establish the displayed scale, weights, exceptional set, or remainder.

\begin{warning}[Former Conjecture 14.4 duplicates the forward report]
\label{ii:warn:defect-duplicate}
The submitted statement that, for every fixed $z>0$, the weighted partition sums
\[
 \sum_{\substack{\pi\in\Pi_m\\1<|\pi|<m}}z^{s(\pi)}2^{-D(\pi)}
\]
admit a full asymptotic expansion whose leading term comes from the two first-defect
shells: one block of size $2$ with the rest singleton, and one block of size $m-1$ with
one singleton is exactly Conjecture 14.3 of the corrected forward-iterate report.
It remains open manuscript mathematics there, with no exact Lean counterpart, but it
is not a second independent inverse-iterate conjecture.  A sharp version would give an
explicit two-spine rate and the first correction term for both forward and prospective
inverse spines.
\end{warning}

\begin{problem}[Nested Lambert phase closure]
\label{ii:prob:nested-lambert}
Put $y_0=y$, $y_{r+1}=I(y_r)$ and, for $0\le r<n$,
\[
 T_r:=\log(1/y_r),\qquad
 \rho_r:=\sqrt{2T_r/L},\qquad
 \vartheta_r:=\rho_r+\beta\pmod{1}.
\]
For fixed $n$ set $q_n:=T_0^{-2^{-n}}$ and $\ell:=\log T_0$.
Determine whether, for every truncation order $N$, there are recursively
computable functions $C_{n,j}$, polynomial in $\ell$ and real analytic and
one-periodic in each of $\vartheta_0,\ldots,\vartheta_{n-1}$, such that
\begin{equation}
 T_n=c_nq_n^{-1}
 +\sum_{j=0}^{N-1}
 C_{n,j}(\ell,\vartheta_0,\ldots,\vartheta_{n-1})q_n^j
 +\BigO_{n,N}\!\left((1+\ell)^{D_{n,N}}q_n^N\right),
 \label{ii:eq:nested-lambert-template}
\end{equation}
where $c_n=(2L)^{1-2^{-n}}$, $D_{n,N}$ is finite, and the remainder
constant is uniform in all phase vectors attained as $y\downarrow0$.
Here ``recursively computable'' is to mean the following explicit
composition-and-truncation procedure.  With the one-step coefficients $h_j$
of \cref{ao:thm:all-orders}, form
\[
 \mathcal I_N(T):=L\rho+L\beta-\log\rho
 -\sum_{j=1}^{N}\frac{h_j(\log\rho,\rho+\beta)}{\rho^j},
 \qquad \rho=\sqrt{2T/L}.
\]
Starting from $\widehat T_0=T_0$, expand
$\mathcal I_N(\widehat T_r)$ in $q_n$, retain all powers below $q_n^N$,
and call the result $\widehat T_{r+1}$.  Prove, uniformly in the attained
phases, that
\[
 T_r-\widehat T_r
 =\BigO_{n,N}\!\left((1+\ell)^{D'_{n,N}}q_n^N\right)
 \qquad(0\le r\le n),
\]
including the phase-error estimates needed before each periodic $h_j$ is
composed with $\sqrt{2\widehat T_r/L}+\beta$.  Finally, determine the first
coefficient in which each $\vartheta_r$ occurs and prove that its dependence
is nonconstant; this is the precise ``periodic ripple is transported''
claim.  If the single lattice $q_n^{\IntegerNumbers}$ is not closed under the
recursion, identify and prove minimality of the larger exponent semigroup.
\end{problem}

The one-step theorem \cref{ao:thm:all-orders} supplies the input to this
problem, and \eqref{ii:eq:iterated-inverse-endpoint} supplies its leading
term.  Neither result by itself justifies the multiscale closure, the phase
substitutions, or the uniform remainder in
\eqref{ii:eq:nested-lambert-template}; the interior Taylor-reversion
experiment tests none of them.

\subsection{Promising formalization route}

A Lean formalization can be modular.
\begin{enumerate}[label=\arabic*.]
 \item Define $Q_m$ and formalize the partition defect and the exact factorization
 \eqref{ii:eq:Q-factorization}.
 \item Prove a sufficiently explicit weighted defect bound for finite set partitions.
 \item Use the existing derivative atlas and bounded-derivative API to formalize the
 two-spine reduction and the fixed-$n$ orbit-spine induction.
 \item Reuse strict monotonicity, convexity, and order-isomorphism modules to prove the
 finite tie-set and dense dominant-amplitude statements.
 \item Finish the forward theorem with the existing power-series radius and analytic-at
 infrastructure.
 \item Formalize the purely algebraic uniqueness of the compositional inverse of a
 formal power series with invertible linear coefficient.
 \item Combine the analytic inverse-function theorem with formal-series convergence to
 prove \cref{ii:thm:formal-reversion} and transport $E_n$ to $Y_n$.
 \item Import the formal quadratic endpoint law and prove the elementary induction for
 \cref{ii:thm:iterated-endpoint}.
\end{enumerate}
The formal-reversion layer should be reusable well beyond the Fabius development.

\subsection{Generalization to other atomic functions}

The proof abstracts beyond the classical dyadic case.  Suppose a smooth increasing map
$f$ has derivatives
\begin{equation}
 f^{(m)}(x)=q^{\binom{m+1}{2}}\Psi(q^m x),
 \label{ii:eq:q-atomic-derivatives}
\end{equation}
for some $q>1$, with bounded $\Psi$, recurrent nonvanishing along a dense set of
$q$-adic orbits, and monotone orbit slopes yielding a unique dominant prefix.  The same
partition calculation replaces $2^{-D(\pi)}$ by $q^{-D(\pi)}$.  Once forward iterates
are controlled, \cref{ii:thm:formal-reversion} automatically transfers every zero-radius
statement to inverse iterates.  This suggests a general theorem for invertible
Rvachev-type distribution functions and multiscale refinable maps.

\section{Summary of the proof mechanism}

The complete argument can be compressed into the following chain.
\begin{enumerate}[label=\arabic*.]
 \item Exact differentiation gives
 $F^{(m)}=Q_m\Theta\comp(2^m\cdot)$ with
 $Q_m=2^{m(m+1)/2}$ and $|\Theta|\leq1$.
 \item In Fa\`a di Bruno's set-partition formula, every non-spine partition has a
 positive quadratic defect $D(\pi)$.
 \item The weighted sum of all non-spine partitions is exponentially small despite
 Bell-number multiplicity.
 \item Hence, after division by $Q_m$, the derivative $G_n^{(m)}$ is asymptotic to a
 finite sum of orbit spines:
 \
 \[
   A_{j,n}B_j^m\Theta(2^m x_j),\qquad 0\leq j<n.
 \]
 \item Away from a finite tie set, monotone orbit dynamics makes one $B_j$ uniquely
 maximal; away from countably many dyadic orbit preimages, binary transitions keep its
 amplitude at least $1/2$ infinitely often.
 \item Therefore $R_{G_n}(x)=0$ on a co-countable dense set $E_n$, and $G_n$ is nowhere
 analytic.
 \item Smooth local inverse Taylor series are unique formal compositional inverses.
 Positive convergence radius is invariant under reversion.
 \item Thus $R_{K_n}(y)=0$ on the co-countable dense set $Y_n=G_n(E_n)$, and no interior
 point of $K_n$ can be analytic.
 \item At the endpoints, $K_n(y)/y^\alpha\to\infty$ for every $\alpha>0$; hence endpoint
 analyticity is impossible.  At the center, the Taylor series is affine but does not
 represent the function.
\end{enumerate}
This proves \cref{ii:thm:inverse-main,ii:thm:inverse-strong} and explains why finite
self-composition cannot regularize either the Fabius function or its inverse.

\appendix

\section{Auxiliary details on the partition defect}
\label{ii:app:defect}

This appendix records two identities useful for sharpening
\cref{ii:prop:defect-decay}.

\begin{proposition}[Exact minimum at fixed block count]
\label{ii:prop:fixed-k-minimum}
Let $m\ge1$ and $1\le k\le m$.  Among partitions of $m$ with $k$ blocks,
\begin{equation}
 \min D(\pi)=(k-1)(m-k).
 \label{ii:eq:fixed-k-minimum}
\end{equation}
Equality holds exactly when one block has size $m-k+1$ and every other block is a
singleton.
\end{proposition}

\begin{proof}
For fixed $m,k$, formula \eqref{ii:eq:defect} shows that minimizing $D$ is equivalent to
maximizing $\sum_i\binom{r_i}{2}$.  Strict convexity of $r\mapsto\binom r2$ transfers
one unit from a smaller nonsingleton block to a larger block and increases the sum.
Iteration leaves one large block and $k-1$ singletons.  Direct substitution gives
\eqref{ii:eq:fixed-k-minimum}.
\end{proof}

\begin{proposition}[First defect shell]
\label{ii:prop:first-shell}
For $m\geq3$, the least positive defect is $m-2$.  It occurs for the two
endpoint forms
\begin{equation}
 (2,1,\ldots,1)
 \qquad\text{and}\qquad
 (m-1,1).
 \label{ii:eq:first-shell}
\end{equation}
which coincide in the boundary case $m=3$ and are distinct for $m\ge4$.
\end{proposition}

\begin{proof}
From \eqref{ii:eq:fixed-k-minimum}, the least positive value over $2\leq k\leq m-1$ is
\[
 \min_{2\leq k\leq m-1}(k-1)(m-k)=m-2,
\]
attained only at the endpoint values $k=2$ and $k=m-1$; these are the same
value when $m=3$.  The equality classification in
\cref{ii:prop:fixed-k-minimum} gives the displayed forms and the stated
distinctness.
\end{proof}

This identifies the candidate leading correction to the spine sum.  At the
near-singleton end, choose one pair and leave all other elements singleton; at the
near-giant end, choose the singleton outside a block of size $m-1$.  Both carry the
same dyadic penalty $2^{-(m-2)}$, but their derivative coefficients and orbit factors
are different.  A refined expansion should sum these shells before estimating the
higher-defect remainder.
